\section*{الفصل \ref{ch:b3:quantum-model-systems} --- ميكانيكا الكم للكيميائيين: الأنظمة النموذجية}

\begin{solution}{exo:b3:quantum-model-systems:1}
$E_1 = h^2/8m_eL^2$. (أ) $L = \qty{1.00}{nm}$: $E_1 = \qty{6.02e-20}{J} =
\qty{0.376}{eV}$؛ $\Delta E_{12} = 3E_1$ و$\lambda = hc/3E_1 =
\qty{1099}{nm}$، في الأشعة تحت الحمراء القريبة. (ب) $L = \qty{0.50}{nm}$: $E_1$ أكبر
بأربع مرات، \qty{2.41e-19}{J} $= \qty{1.50}{eV}$، و$\lambda =
\qty{275}{nm}$، في الأشعة فوق البنفسجية.
\end{solution}

\begin{solution}{exo:b3:quantum-model-systems:2}
$n_x^2 + n_y^2 + n_z^2 = 3$ (1,1,1): $g = 1$؛ $6$ (2,1,1 وتبديلاتها):
$g = 3$؛ $9$ (2,2,1): $g = 3$؛ $11$ (3,1,1): $g = 3$؛ $12$ (2,2,2): $g = 1$.
\end{solution}

\begin{solution}{exo:b3:quantum-model-systems:3}
$[\hat x^2,\hat p]f = -\iu\hbar x^2f' + \iu\hbar(x^2f)' = 2\iu\hbar xf$، ومنه
$[\hat x^2,\hat p] = 2\iu\hbar\hat x$. و$[\hat p^2,\hat x]f = -\hbar^2[(xf)'' -
xf''] = -2\hbar^2f'$، ومنه $[\hat p^2,\hat x] = -2\iu\hbar\hat p$.
\end{solution}

\begin{solution}{exo:b3:quantum-model-systems:4}
$E_J/hc = BJ(J+1)$: 0 و21.19 و63.56 و\qty{127.12}{cm^{-1}}، مع $g = 1$
و3 و5 و7. $I = h/(8\pi^2cB) = \qty{2.643e-47}{kg.m^2}$ (مع $c$ بوحدة
\unit{cm/s}).
\end{solution}

\begin{solution}{exo:b3:quantum-model-systems:5}
$|\psi_n|^2$ متناظرة بالنسبة إلى $L/2$، ومنه $\langle x\rangle = L/2$. وباستعمال
$\sin^2u = \frac12(1 - \cos2u)$ ومكاملتين بالتجزئة،
$\langle x^2\rangle = L^2\big(\frac13 - \frac{1}{2n^2\pi^2}\big)$. ومن أجل $n = 1$،
$\langle x^2\rangle = 0.283L^2$، وهي أقل من القيمة الكلاسيكية $L^2/3$: فالحالة
الأساسية تتكدس في الوسط. وكلما كبر $n$ استُعيدت القيمة الكلاسيكية.
\end{solution}

\begin{solution}{exo:b3:quantum-model-systems:6}
$1\,\unit{cm^{-1}} = \qty{0.011963}{kJ/mol}$. \omterm{def:b3:quantum-model-systems:oscillator}{طاقة النقطة الصفرية} (\ce{H2}) $= 2200.6 \times
0.011963 = \qty{26.33}{kJ/mol}$، \omterm{def:b3:quantum-model-systems:oscillator}{وطاقة النقطة الصفرية} (\ce{D2}) $= \qty{18.63}{kJ/mol}$؛ والفرق
\qty{7.69}{kJ/mol}. \omterm{def:b3:quantum-model-systems:oscillator}{ثابت القوة} نفسه، ومنه $\tilde\omega \propto \mu^{-1/2}$
والنسبة المتوقعة $\sqrt{\mu_D/\mu_H} = 1.4137$ (الكتلتان الذريتان
2.01410 و1.00783)؛ والنسبة المقيسة 1.4127، أي في حدود 0.07\,\% (
والفرق الصغير مصدره الإلكترونات، التي لا تتبع النوى
تمامًا).
\end{solution}

\begin{solution}{exo:b3:quantum-model-systems:7}
$\int_0^Lx^2(L - x)^2\,\dd x = L^5/30$، ومنه $N = \sqrt{30/L^5}$. عندئذ
\[
\langle\psi_1|\phi\rangle = \sqrt{\tfrac2L}\sqrt{\tfrac{30}{L^5}}\int_0^Lx(L -
x)\sin\tfrac{\pi x}{L}\,\dd x = \frac{\sqrt{60}}{L^3}\cdot\frac{4L^3}{\pi^3} =
\frac{4\sqrt{60}}{\pi^3} = 0.99928 .
\]
القطع المكافئ حالة أساسية بنسبة 99.86\,\% (مربع التداخل).
\end{solution}

\begin{solution}{exo:b3:quantum-model-systems:8}
$\langle f|\hat pg\rangle = -\iu\hbar\int f^*g'\,\dd x = -\iu\hbar[f^*g] +
\iu\hbar\int f^{*\prime}g\,\dd x = \int(-\iu\hbar f')^*g\,\dd x = \langle\hat
pf|g\rangle$، إذ ينعدم ما بين المعقوفين عند الطرفين. ومن أجل $\hat x\hat p$، وباستعمال أن كلًا من $\hat x$ و$\hat p$ هرميتي:
$\langle f|\hat x\hat pg\rangle = \langle\hat xf|\hat pg\rangle = \langle\hat p\hat
xf|g\rangle$، و$\hat p\hat x = \hat x\hat p - \iu\hbar \ne \hat x\hat p$: فهو ليس
\omterm{def:b3:quantum-model-systems:hermitian}{هرميتيًا} (أما صيغته المتناظرة $\frac12(\hat x\hat p + \hat p\hat x)$ فهرميتية).
\end{solution}

\begin{solution}{exo:b3:quantum-model-systems:9}
$L = 6 \times 139.7 = \qty{838}{pm}$، $N = 6$: $\lambda = 8m_ecL^2/(7h) =
\qty{331}{nm}$. يقع الامتصاص في الأشعة فوق البنفسجية: فالهكساترايين
عديم اللون.
\end{solution}

\begin{solution}{exo:b3:quantum-model-systems:10}
$\kappa = \sqrt{2m(V_0 - E)}/\hbar$ مع $V_0 - E = \qty{0.20}{eV}$:
$\kappa_H = \qty{9.82e10}{m^{-1}}$، $\kappa_D = \qty{1.388e11}{m^{-1}}$. 
وتتلاشى المعاملات القبلية في النسبة: $T_H/T_D = \eu^{2a(\kappa_D - \kappa_H)} =
\eu^{4.06} = 58$، وهي القيمة الدقيقة: فالحاجز سميك لكليهما
($\kappa a = 4.9$ و6.9).
\end{solution}

\begin{solution}{exo:b3:quantum-model-systems:11}
المستويات $E_m = m^2\hbar^2/2m_eR^2$: يتسع $m = 0$ لإلكترونين، و$m = \pm1$
لأربعة. أعلى مستوى مملوء هو $|m| = 1$، وأدنى مستوى فارغ $|m| = 2$: $\Delta E =
3\hbar^2/2m_eR^2 = \qty{9.38e-19}{J}$، $\lambda = \qty{212}{nm}$، في الأشعة
فوق البنفسجية، حيث يقع فعلًا الامتصاص القوي للبنزين.
\end{solution}

\begin{solution}{exo:b3:quantum-model-systems:12}
$\langle r\rangle = \frac{4\pi}{\pi a^3}\int_0^\infty r^3\eu^{-2r/a}\dd r =
\frac{4}{a^3}\cdot\frac{3!}{(2/a)^4} = \frac32a$، أي \qty{79.4}{pm} من أجل $a = a_0$.
$\langle 1/r\rangle = \frac{4}{a^3}\cdot\frac{1!}{(2/a)^2} = \frac1a$، ومنه
$\langle V\rangle = -e^2/(4\pi\varepsilon_0a)$. وبما أن $E_1 =
-e^2/(8\pi\varepsilon_0a)$ (انطلاقًا من $a = 4\pi\varepsilon_0\hbar^2/\mu e^2$)،
فإن $\langle V\rangle = 2E_1$، ومتوسط الطاقة الحركية $-E_1$.
\end{solution}

\begin{solution}{pb:b3:quantum-model-systems:1}
\textbf{1.} تقدّم كل ذرة من الذرات $j$ مدارًا p واحدًا؛ وتقدّم ذرات الكربون $j - 1$ 
وإحدى ذرتي النيتروجين إلكترونًا واحدًا لكل منها، وتقدّم ذرة النيتروجين الأخرى زوجها الحر،
لأن شحنة الكاتيون موزعة على السلسلة: $N = j + 1$، أي 6
و8 و10.
\textbf{2.} ستة إلكترونات تملأ $n = 1$ و2 و3: HOMO هو $n = 3$، وLUMO هو $n = 4$.
\textbf{3.} $\Delta E = E_{N/2+1} - E_{N/2} = (N+1)h^2/8mL^2$ و$\lambda =
hc/\Delta E = 8mcL^2/h(N+1)$.
\textbf{4.} $L_0 = 4l$ و$6l$ و$8l$: 559 و838 و\qty{1118}{pm}.
\textbf{5.} 147 و257 و\qty{374}{nm}.
\textbf{6.} كلها أقصر بكثير (بعوامل 3.6 و2.3 و1.9): صندوق النموذج
صغير جدًا، لأن $\lambda \propto L^2$؛ فالإلكترونات تتحرك على مسافة أطول
من السلسلة الممتدة بين ذرتي النيتروجين.
\textbf{7.} بأخذ المبدأ في المركز، $\psi_n \propto \cos(n\pi u/L)$ من أجل
$n$ فردي و$\sin(n\pi u/L)$ من أجل $n$ زوجي ($u = x - L/2$): أي دوال زوجية وفردية
في $u$.
\textbf{8.} إذا كان $n + n'$ زوجيًا، كان للدالتين $\psi_n$ و$\psi_{n'}$ الزوجية نفسها؛
فجداؤهما مضروبًا في $u$ دالة فردية تكاملها معدوم على فترة
متناظرة.
\textbf{9.} للمستويين HOMO $N/2$ وLUMO $N/2 + 1$ مجموع $n + n' = N + 1$، وهو فردي: فالانتقال
مسموح دائمًا.
\textbf{10.} الصبغة الثانية: HOMO هو 4، وLUMO هو 5. $4 \to 6$: المجموع زوجي، ممنوع.
$3 \to 5$: المجموع زوجي، ممنوع.
\textbf{11.} الانتقال المسموح التالي انطلاقًا من $n = 4$ هو $4 \to 7$، 
و$\Delta E$ فيه أكبر بالعامل $(49 - 16)/(25 - 16) = 3.67$: أي عند $\lambda/3.67$،
\qty{164}{nm} للصبغة الثانية، في عمق الأشعة فوق البنفسجية.
\textbf{12.} إنها لا تتعلق إلا بتناظر الكمون بالنسبة إلى مركز
السلسلة، وهو تناظر تشترك فيه الصبغات المتناظرة الحقيقية.
\textbf{13.} $E = hc/\lambda = \qty{3.294e-19}{J} = \qty{2.06}{eV}$.
\textbf{14.} $hc\tilde\omega_e = \qty{0.371}{eV}$.
\textbf{15.} $2hcB = \qty{4.21e-22}{J} = \qty{2.63}{meV}$.
\textbf{16.} $kT = \qty{25.7}{meV}$.
\textbf{17.} المستويات الدورانية وحدها ($2hcB \ll kT$)؛ أما الإثارات الاهتزازية ($15kT$) 
والإلكترونية ($80kT$) فغير مأهولة عمليًا.
\textbf{18.} الإلكترونية: المرئي وفوق البنفسجي؛ والاهتزازية: تحت الأحمر؛
والدورانية: الأمواج الميكروية (وتحت الأحمر البعيد).
\textbf{19.} $L = \sqrt{\lambda h(N+1)/8mc}$ مع $N = 8$: $L = \qty{1283}{pm}$.
\textbf{20.} $\delta = (1283 - 838)/2 = \qty{222}{pm}$.
\textbf{21.} $L = 559 + 445 = \qty{1003}{pm}$، $\lambda = \qty{474}{nm}$، أي 9.5\,\%
دون 524 نانومترًا.
\textbf{22.} $L = 1118 + 445 = \qty{1562}{pm}$، $\lambda = \qty{732}{nm}$، أي 2.9\,\%
فوق 711 نانومترًا.
\textbf{23.} إلى الحلقتين العطريتين اللتين تحملان ذرتي النيتروجين: فنظام $\pi$
لا ينتهي عند ذرتي النيتروجين، والصندوق الفعلي يضم
جزءًا من كل حلقة.
\textbf{24.} $\delta/l = 222/139.7 = 1.6$: \textbf{يمتد الصندوق نحو
\qty{222}{pm}، أي رابطة ونصفًا، وراء كل ذرة نيتروجين}، وبهذا الطول الواحد
يعيد النموذج الألوان الثلاثة في حدود 10\,\%.
\end{solution}
